## topex / tas_mean nonlinearity + tas_mean ICC check (2026-08-20)

Question: does Elastic Net's linear-only assumption explain topex's (and tas_mean's) sign
instability, as a mechanism separate from / alongside the spatial-confounding (ICC) account
already in the chapter? Checked directly, no retraining -- raw data only
(`data/processed/environmental/plot_environmental_features.parquet`, 4survey, n=71,766).

### Script (reproducible, no saved-model dependency)

```python
import pandas as pd
import numpy as np
from scipy import stats

df = pd.read_parquet('data/processed/environmental/plot_environmental_features.parquet')
sub = df.dropna(subset=['mean_cr_residual_4survey','tas_mean_4survey','topex','cpmt']).copy()

# 1. coarseness
print('tas_mean_4survey unique values:', df['tas_mean_4survey'].nunique(), '/', df['tas_mean_4survey'].notna().sum())
print('topex unique values:', df['topex'].nunique(), '/', df['topex'].notna().sum())

# 2. bivariate linearity check (Pearson vs Spearman)
for col in ['tas_mean_4survey','topex']:
    pear = stats.pearsonr(sub[col], sub['mean_cr_residual_4survey'])
    spear = stats.spearmanr(sub[col], sub['mean_cr_residual_4survey'])
    print(col, 'pearson', pear[0], 'spearman', spear[0])

# 3. binned means (deciles) -- catches non-monotonic shape Spearman misses
for col in ['tas_mean_4survey', 'topex']:
    sub['bin'] = pd.qcut(sub[col], 10, duplicates='drop')
    print(sub.groupby('bin', observed=True)['mean_cr_residual_4survey'].agg(['mean','count']))

# 4. between-compartment variance share (ICC), same method as topex's existing 74% figure
def icc(frame, col):
    total_var = frame[col].var(ddof=0)
    comp_means = frame.groupby('cpmt')[col].transform('mean')
    return comp_means.var(ddof=0) / total_var

for col in ['tas_mean_4survey','topex']:
    print(col, 'between-compartment variance share', icc(sub, col))
```

### Results

**Coarseness**: tas_mean_4survey = 149 distinct values / 71,766 plots (1km HadUK-Grid).
topex = 16,839 distinct values / 71,766 (much finer).

**Bivariate Pearson vs Spearman** (crude linearity check -- does NOT catch non-monotonic shape):
- tas_mean_4survey: pearson r=0.2794, spearman rho=0.2733 -- nearly identical, no flag.
- topex: pearson r=0.0803, spearman rho=0.0930 -- nearly identical, no flag.

**Binned means** (10 deciles, x=variable value, y=mean of mean_cr_residual_4survey in that bin):

topex deciles (low to high): -2.28, -0.43, -0.12, +0.35, -0.95, +1.34, +1.36, +1.39, +0.16, -0.92.
Clearly non-monotonic -- down/up/down/up-plateau/down. Not a line, not even a single monotonic
curve. This is why the crude Pearson/Spearman check above missed it: a wave that nets out near
zero rank-correlation can still hide large local swings a linear coefficient cannot represent
with one slope.

tas_mean deciles (coldest to warmest): -3.95, -1.58, -0.85, +0.48, -0.28, +0.77, +0.99, +2.04,
+2.53, -0.09. Mostly rising, but reverses sharply at the warmest decile.

**Between-compartment variance share (ICC)**, same method as the topex 74% figure already in the
chapter (recomputed here too as a consistency check: 0.7428, matches):
- topex: 0.7428 (matches chapter's existing 74%)
- tas_mean_4survey: **0.8933** -- higher than topex, i.e. tas_mean is even more spatially
  confounded with compartment structure than topex is.

### Read

Two separate, complementary (not competing) mechanisms for topex's EN instability:
1. Spatial confounding (already in chapter): topex mostly varies between compartments, EN has no
   random intercept to absorb that, LMM does.
2. NEW, this check: topex's own relationship with the target is non-monotonic. A single linear EN
   coefficient is an unstable summary of a wavy shape -- which fold's data samples more of the
   "up" vs "down" segments can flip the fitted sign. LMM isn't immune to this either (it's still
   linear in fixed effects) but its random intercept absorbs enough compartment-level structure
   that the residual relationship it's fitting may be flatter than EN's -- not directly tested
   here, flagged as a further check if wanted.

tas_mean: both mechanisms apply, and more strongly on the ICC front (89% vs topex's 74%) --
consistent with tas_mean also being EN-unstable (matches the chapter's own numbers: VIF 3.44,
named as one of the unstable ones) and gives it an evidenced account, not just an assumed one.

### Addendum (same day) -- soilgrids_ph + chelsa_bio12_precip_mm, completing the check

Ran the same two diagnostics (binned deciles, ICC) on the remaining two named variables in the
paragraph: soilgrids_ph (unstable) and chelsa_bio12_precip_mm (stable, the actual CHELSA/
1981-2010 variable). This changes the reading.

**soilgrids_ph**: 13 distinct values total (coarser than tas_mean's 149). ICC = 0.774. Binned
means (5 bins, qcut collapsed from 10 due to few unique values): -3.18, +0.16, +1.17, +1.70,
-0.30 -- rises then reverses at the top bin, same non-monotonic shape as topex/tas_mean.

**chelsa_bio12_precip_mm** (stable): ICC = **0.973** -- higher than topex (0.743), tas_mean
(0.893), AND soilgrids_ph (0.774). Binned means: 2.35, 2.21, 1.70, 1.37, 0.34, -0.45, -1.18,
-1.21, -0.83, -4.57 -- close to monotonic decreasing, only one small local wiggle. Pearson
-0.298, Spearman -0.364 -- clearly the strongest bivariate signal of all four variables checked.

**Revised read**: ICC alone does not separate stable from unstable -- the highest-ICC variable of
the four is stable. What differs: chelsa_bio12_precip_mm has a STRONG, near-linear relationship
with the target; the three unstable variables (topex, tas_mean, soilgrids_ph) all have WEAK
bivariate signal (Spearman 0.09-0.27, vs chelsa_bio12's 0.36) AND a non-monotonic/wavy shape.
Weak + nonlinear -> unstable, regardless of ICC. Strong + near-linear -> stable, even under
higher spatial confounding than any of the unstable ones. Matches the chapter's own existing
slope_degrees-vs-topex point (weaker signal needs less confounding pressure to flip sign) --
this generalises it to all four variables checked, with numbers.

### Second addendum (2026-08-20, later) -- within-compartment decomposition supersedes the
### "nonlinearity as a separate mechanism" framing above

User's catch: LMM is also a LINEAR model (one coefficient per variable, no curve fit) -- if
topex's raw non-monotonic shape were the reason EN is unstable, LMM should be unstable too. It
isn't. The only structural difference between LMM and EN is LMM's compartment random intercept.
So the "wave" found above is not necessarily a mechanism independent of spatial confounding --
it may just be what spatial confounding LOOKS like in a raw, pooled bivariate plot.

Tested directly: subtract each compartment's own mean from both the variable and
mean_cr_residual_4survey (`x - groupby('cpmt').transform('mean')`, both sides), then redo the
decile/correlation check on these WITHIN-compartment residuals. This is mechanically close to
what a compartment random intercept absorbs.

| Variable | Raw Spearman | Within-compartment Spearman | Within-compartment shape |
|---|---:|---:|---|
| topex | 0.093 | **0.210** (stronger) | Close to monotonic, one small dip -- the wave mostly disappears |
| tas_mean_4survey | 0.273 | **0.082** (much weaker) | Roughly monotonic, but little signal left |
| soilgrids_ph | 0.271 | **0.046** (near zero) | No clear shape, signal essentially gone |
| chelsa_bio12_precip_mm (stable) | -0.364 | -0.114 (weaker, but real) | Still clean, monotonically decreasing |

### Revised read (supersedes the "two independent mechanisms" framing above)

Not two separate, complementary mechanisms (confounding + nonlinearity). One mechanism, three
different flavours of it:

- **topex**: real within-compartment signal EXISTS and is fairly clean/monotonic (0.210,
  stronger than the raw 0.093) -- but EN fits the raw, pooled-across-compartments relationship,
  which looks weaker and wavy because different compartments sit at different points on both
  topex and the residual. LMM's random intercept removes this pooling artifact and sees the
  cleaner within-compartment signal directly -- consistent with LMM's coefficient being strong
  and stable.
- **tas_mean, soilgrids_ph**: within-compartment signal nearly vanishes (0.082, 0.046) --
  almost all of their raw bivariate correlation IS the between-compartment difference, not a
  real, generalisable within-plot relationship. There is little genuine signal for any model to
  stably estimate a coefficient from in the first place -- EN's instability here reflects real
  signal scarcity, not (only) a shape problem.
- **chelsa_bio12_precip_mm** (stable): retains real, clean, monotonic within-compartment signal
  (-0.114) on top of very high between-compartment variance (0.97 ICC) -- enough genuine local
  signal to anchor a stable EN coefficient even under heavy confounding.

This is a cleaner, more mechanistically unified account than "confounding OR nonlinearity" --
whether a variable's raw signal survives a within-compartment decomposition (topex,
chelsa_bio12_precip_mm) or evaporates (tas_mean, soilgrids_ph) is what predicts EN stability,
not the raw shape or the raw correlation size alone.

### Caveat

Bivariate/ICC only -- not the actual multivariate Set3/Set4 EN fit (which includes topex/
tas_mean alongside every other predictor at once). Real, direct evidence, but complementary to,
not proof of, the specific coefficient behaviour inside the full model. No SHAP-dependence
(partial, multivariate-adjusted) version of this check has been run.
